\begingroup
\setlength{\parskip}{.3\baselineskip}
La posición de un sabor quark se especifica mediante cuatro datos simultáneos: rama, generación, soporte celular y familia de rutas orientadas. Las tres generaciones de una misma rama comparten el soporte del atlas; el clasificador fino de generación conserva la distinción física que todavía debe realizarse como autovector y sección de masa.
\begingroup
\setlength{\tabcolsep}{4.5pt}
\renewcommand{\arraystretch}{1.28}
\begin{longtable}{@{}>{\raggedright\arraybackslash}p{.12\textwidth}>{\raggedright\arraybackslash}p{.17\textwidth}>{\raggedright\arraybackslash}p{.19\textwidth}>{\raggedright\arraybackslash}p{.46\textwidth}@{}}
\caption{Clasificación generacional de los seis sabores quark.}\label{tab:masas-posicion-quarks-seis-sabores}\\
\toprule
\textbf{Sabor} & \textbf{Generación} & \textbf{Rama} & \textbf{Clasificador de rama y generación} \\
\midrule
\endfirsthead
\toprule
\textbf{Sabor} & \textbf{Generación} & \textbf{Rama} & \textbf{Clasificador de rama y generación} \\
\midrule
\endhead
\midrule
\multicolumn{4}{r}{\footnotesize Continúa en la página siguiente.}\\
\endfoot
\bottomrule
\endlastfoot
$u/\bar u$ & primera & superior & $(\operatorname{up},1,\chi_{\rm fina})$; par nominal de ruta $(21,30)$ sobre el eje $\mathrm{N\!-\!S}$ \\
$d/\bar d$ & primera & inferior & $(\operatorname{down},1,\chi_{\rm fina})$; par nominal de ruta $(20,32)$ sobre el eje $\mathrm{E\!-\!O}$ \\
$c/\bar c$ & segunda & superior & $(\operatorname{up},2,\chi_{\rm fina})$; par nominal de ruta $(21,29)$ sobre el eje $\mathrm{E\!-\!O}$ \\
$s/\bar s$ & segunda & inferior & $(\operatorname{down},2,\chi_{\rm fina})$; par nominal de ruta $(17,27)$ sobre el eje $\mathrm{N\!-\!S}$ \\
$t/\bar t$ & tercera & superior & $(\operatorname{up},3,\chi_{\rm fina})$; par nominal de ruta $(21,26)$ sobre el eje $\mathrm{N\!-\!S}$ \\
$b/\bar b$ & tercera & inferior & $(\operatorname{down},3,\chi_{\rm fina})$; par nominal de ruta $(21,30)$ sobre el eje $\mathrm{E\!-\!O}$ \\
\end{longtable}
\endgroup

\Needspace{18\baselineskip}
\noindent\textbf{Soporte estructural de la rama superior.}\par
\noindent\textit{Multisecciones.} \(\mathcal B_{u,\mathrm{G1}}\), \(\mathcal B_{u,\mathrm{G3}}\).\par
\noindent\textit{Profundidades y cardinales.} $\mathrm{G1}$, $\mathrm{G3}$; 16 celdas y 14 familias de ruta distintas.\par
\noindent\textit{Órbita cromática y orientaciones.} $\mathcal C=\{B,G,R\}$ y $\mathcal O=\{(-1,-1),(-1,1),(1,-1),(1,1)\}$.\par
\noindent\textit{Celdas y familias orientadas $B_1$.} Cada cadena $k\leftrightarrow(r,c)\leftrightarrow\Gamma_{r,c}^{(h,v)}$ conserva el índice celular, su coordenada en la carta $9\times9$ y los dos signos de orientación del registro de ruta enriquecido.\par
\noindent Fila $r=2$: $11\leftrightarrow(2,2)\leftrightarrow\Gamma_{2,2}^{(-,+)}$, $13\leftrightarrow(2,4)\leftrightarrow\Gamma_{2,4}^{(+,-)}$, $15\leftrightarrow(2,6)\leftrightarrow\Gamma_{2,6}^{(-,+)}$, $17\leftrightarrow(2,8)\leftrightarrow\Gamma_{2,8}^{(-,-)}$.\par
\noindent Fila $r=4$: $29\leftrightarrow(4,2)\leftrightarrow\Gamma_{4,2}^{(+,-)}$, $31\leftrightarrow(4,4)\leftrightarrow\Gamma_{4,4}^{(-,+)}$, $33\leftrightarrow(4,6)\leftrightarrow\Gamma_{4,6}^{(+,+)}$, $35\leftrightarrow(4,8)\leftrightarrow\Gamma_{4,8}^{(-,-)}$.\par
\noindent Fila $r=6$: $47\leftrightarrow(6,2)\leftrightarrow\Gamma_{6,2}^{(-,+)}$, $49\leftrightarrow(6,4)\leftrightarrow\Gamma_{6,4}^{(-,-)}$, $51\leftrightarrow(6,6)\leftrightarrow\Gamma_{6,6}^{(+,+)}$, $53\leftrightarrow(6,8)\leftrightarrow\Gamma_{6,8}^{(+,+)}$.\par
\noindent Fila $r=8$: $65\leftrightarrow(8,2)\leftrightarrow\Gamma_{8,2}^{(+,+)}$, $67\leftrightarrow(8,4)\leftrightarrow\Gamma_{8,4}^{(+,+)}$, $69\leftrightarrow(8,6)\leftrightarrow\Gamma_{8,6}^{(-,-)}$, $71\leftrightarrow(8,8)\leftrightarrow\Gamma_{8,8}^{(+,+)}$.\par
\noindent\textit{Morfismo de realización cerrado.} Para $g\in\{1,2,3\}$, la rama, la generación, la órbita de color y la firma fina determinan
\[\chi_{\rm fina}(\operatorname{up},g,\mathcal C)\longrightarrow\Gamma_{\rm enr}\longrightarrow\bigl(\operatorname{quot}_{\mathcal C},S(B_D,B_\Omega)\bigr).\]
Esta flecha selecciona el autovector físico y publica después la coordenada escalar en $\mathcal L_M$ sin usar el valor de contraste.\par
\medskip
\Needspace{18\baselineskip}
\noindent\textbf{Soporte estructural de la rama inferior.}\par
\noindent\textit{Multisecciones.} \(\mathcal B_{d,\mathrm{G1}}\), \(\mathcal B_{d,\mathrm{G2}}\), \(\mathcal B_{d,\mathrm{G3}}\), \(\mathcal B_{d,\mathrm{G4}}\).\par
\noindent\textit{Profundidades y cardinales.} $\mathrm{G1}$, $\mathrm{G2}$, $\mathrm{G3}$, $\mathrm{G4}$; 20 celdas y 18 familias de ruta distintas.\par
\noindent\textit{Órbita cromática y orientaciones.} $\mathcal C=\{B,G,R\}$ y $\mathcal O=\{(-1,-1),(-1,1),(1,-1),(1,1)\}$.\par
\noindent\textit{Celdas y familias orientadas $B_1$.} Cada cadena $k\leftrightarrow(r,c)\leftrightarrow\Gamma_{r,c}^{(h,v)}$ conserva el índice celular, su coordenada en la carta $9\times9$ y los dos signos de orientación del registro de ruta enriquecido.\par
\noindent Fila $r=2$: $10\leftrightarrow(2,1)\leftrightarrow\Gamma_{2,1}^{(+,+)}$, $12\leftrightarrow(2,3)\leftrightarrow\Gamma_{2,3}^{(-,-)}$, $14\leftrightarrow(2,5)\leftrightarrow\Gamma_{2,5}^{(+,+)}$, $16\leftrightarrow(2,7)\leftrightarrow\Gamma_{2,7}^{(-,+)}$, $18\leftrightarrow(2,9)\leftrightarrow\Gamma_{2,9}^{(-,+)}$.\par
\noindent Fila $r=4$: $28\leftrightarrow(4,1)\leftrightarrow\Gamma_{4,1}^{(-,-)}$, $30\leftrightarrow(4,3)\leftrightarrow\Gamma_{4,3}^{(+,+)}$, $32\leftrightarrow(4,5)\leftrightarrow\Gamma_{4,5}^{(+,-)}$, $34\leftrightarrow(4,7)\leftrightarrow\Gamma_{4,7}^{(-,+)}$, $36\leftrightarrow(4,9)\leftrightarrow\Gamma_{4,9}^{(-,+)}$.\par
\noindent Fila $r=6$: $46\leftrightarrow(6,1)\leftrightarrow\Gamma_{6,1}^{(-,-)}$, $48\leftrightarrow(6,3)\leftrightarrow\Gamma_{6,3}^{(-,+)}$, $50\leftrightarrow(6,5)\leftrightarrow\Gamma_{6,5}^{(-,+)}$, $52\leftrightarrow(6,7)\leftrightarrow\Gamma_{6,7}^{(-,+)}$, $54\leftrightarrow(6,9)\leftrightarrow\Gamma_{6,9}^{(+,-)}$.\par
\noindent Fila $r=8$: $64\leftrightarrow(8,1)\leftrightarrow\Gamma_{8,1}^{(+,-)}$, $66\leftrightarrow(8,3)\leftrightarrow\Gamma_{8,3}^{(-,+)}$, $68\leftrightarrow(8,5)\leftrightarrow\Gamma_{8,5}^{(-,+)}$, $70\leftrightarrow(8,7)\leftrightarrow\Gamma_{8,7}^{(+,-)}$, $72\leftrightarrow(8,9)\leftrightarrow\Gamma_{8,9}^{(-,+)}$.\par
\noindent\textit{Morfismo de realización cerrado.} Para $g\in\{1,2,3\}$, la rama, la generación, la órbita de color y la firma fina determinan
\[\chi_{\rm fina}(\operatorname{down},g,\mathcal C)\longrightarrow\Gamma_{\rm enr}\longrightarrow\bigl(\operatorname{quot}_{\mathcal C},S(B_D,B_\Omega)\bigr).\]
Esta flecha selecciona el autovector físico y publica después la coordenada escalar en $\mathcal L_M$ sin usar el valor de contraste.\par
\medskip
\endgroup
